% !PRINTED 171
\weberpage{171}
Finally we may also inquire after fractional invariants. If
\[
\frac{F(x)}{F_1(x)}
\]
is such a fractional invariant, we first assume that the two integral rational functions $F(x), F_1(x)$ of the $n$ variables $x$ are freed of common divisors (Bd.\,I, \S.\,51).

If we restrict ourselves for the present to absolute invariants, and if accordingly
\[
\frac{F(x)}{F_1(x)} = \frac{F[A(x)]}{F_1[A(x)]},
\]
then, since the $x$ are here regarded as independent variables, numerator and denominator have no common divisor either on the right or on the left, and it follows that:
\[
F[A(x)] = \alpha F(x), \quad F_1[A(x)] = \alpha F_1(x),
\]
in which $\alpha$ is a constant factor which, as we have seen earlier, is a root of unity. Numerator and denominator of a fractional invariant must therefore themselves be invariants, if only relative ones.

But if we are not seeking precisely the simplest representation, we can also represent absolute fractional invariants as quotients of absolute integral invariants. We need only extend the fraction by a suitable power of the denominator, that is, if the $\alpha$ are $e$th roots of unity, to put
\[
\frac{F(x)}{F_1(x)} = \frac{F(x) F_1(x)^{e-1}}{[F_1(x)]^e}
\]
Accordingly we can also represent all relative invariants with a definite system of factors $\alpha, \beta, \ldots$ as the product of one of them with absolute invariants, which however may be fractional.

\begin{center}
\S.\,43.

\gesperrt{The Form Problem.}
\end{center}

In the preceding section it was proved that for a finite group of linear substitutions there is a finite number of independent invariants. If $n$ is the dimension of the linear
%CONT>>
